\documentclass[11pt]{article}
\usepackage[margin=1in]{geometry}
\usepackage{amsmath,amssymb,amsthm}

\newtheorem{theorem}{Theorem}
\newtheorem{lemma}{Lemma}
\theoremstyle{remark}
\newtheorem{remark}{Remark}

\newcommand{\ad}{\operatorname{ad}}
\newcommand{\im}{\operatorname{im}}
\newcommand{\spn}{\operatorname{span}}
\newcommand{\one}{\mathbf{1}}

\title{Local Order of Strang Splitting for Finite-State Markov Generators:\\
a Short Proof}
\author{}
\date{}

\begin{document}
\maketitle

\section{Statement}

For square matrices $B,C$ of the same size put
$$
A=B+C,\qquad S(h)=e^{hB/2}e^{hC}e^{hB/2},\qquad P(h)=e^{hA}.
$$
The splitting is \emph{exact} if $S(h)=P(h)$ for all real $h$. Otherwise
$p(B,C)$ denotes the order of vanishing of $S(h)-P(h)$ at $h=0$:
$$
S(h)-P(h)=h^{p}K+O(h^{p+1}),\qquad K\neq 0 .
$$
A \emph{Markov generator} is a real matrix with nonnegative off-diagonal
entries and zero row sums. For $n\ge 2$ let
$$
M_n=\sup\{\,p(B,C): B,C \text{ are } n\times n \text{ Markov generators
with nonexact splitting}\,\}.
$$

\begin{theorem}\label{thm:main}
For every $n\ge 2$,
$$
2\left\lfloor \frac n2\right\rfloor+1\;\le\; M_n\;\le\;\max\{3,\,(n-1)(n-2)+1\}.
$$
The lower bound is attained by generators whose off-diagonal entries are
all strictly positive. In particular $M_2=M_3=3$.
\end{theorem}

Throughout, $[U,V]=UV-VU$, $\ad_U V=[U,V]$, and we fix the notation
$$
\boxed{\;D=B+2C,\qquad T=\ad_B,\qquad Z=\ad_D.\;}
$$
We use $e^{U}Ve^{-U}=e^{\ad_U}V$ freely. Lemmas~\ref{lem:parity}--\ref{lem:arb}
hold for arbitrary complex square matrices.

\section{The order criterion}

\begin{lemma}[Parity]\label{lem:parity}
If the splitting is nonexact, then $p(B,C)$ is odd.
\end{lemma}

\begin{proof}
For small $h$ the principal logarithm $\Lambda(h)=\log S(h)$ is analytic,
and $S(-h)=S(h)^{-1}$ gives $\Lambda(-h)=-\Lambda(h)$. Since
$S(h)=I+hA+O(h^2)$, we have $\Lambda(h)=hA+O(h^2)$. If $\Lambda(h)=hA$
identically near $0$, then $S=P$ near $0$ and hence everywhere by
analyticity. Otherwise $\Lambda(h)=hA+h^{q}K+O(h^{q+1})$ with $K\ne0$,
and $q$ is odd because $\Lambda$ is odd. Both $\Lambda(h)$ and $hA$ are
$O(h)$, so
$$
S(h)-P(h)=e^{\Lambda(h)}-e^{hA}=h^{q}K+O(h^{q+1}),
$$
and $p=q$.
\end{proof}

\begin{lemma}[Order criterion]\label{lem:crit}
The splitting is exact if and only if $Z\,T^{s}D=0$ for every $s\ge1$,
that is, if and only if $D$ commutes with $\ad_B^{s}D$ for all $s\ge 1$.
Otherwise
$$
p(B,C)=2+\min\{\,s\ge 1: Z\,T^{s}D\neq 0\,\}.
$$
In particular $p\ge 3$, and $p>3$ if and only if $[D,[B,D]]=0$.
\end{lemma}

\begin{proof}
\emph{Step 1: a defect function.}
Let $E=S-P$ and $L=S'S^{-1}-A$. Since $P'=AP$,
\begin{equation}\label{eq:LE}
L=(E'-AE)\,S^{-1}.
\end{equation}
Differentiating the three factors of $S$,
$$
S'S^{-1}=\tfrac12 B+e^{hT/2}C+\tfrac12\,e^{hT/2}e^{h\ad_C}B .
$$
Put $R=e^{-hT/2}L$. Using $e^{-hT/2}B=B$,
$$
R(h)=\tfrac12\bigl(e^{h\ad_C}-I\bigr)B-\bigl(e^{-hT/2}-I\bigr)C,
\qquad
R'(h)=\tfrac12\bigl(e^{-hT/2}-e^{h\ad_C}\bigr)Q,
\quad Q=[B,C].
$$

\emph{Step 2: the order of $R'$ is $p-2$.}
If $E\equiv0$ then $L\equiv0$ by \eqref{eq:LE}, so $R'\equiv0$.
Conversely, if $R'\equiv0$ then $R\equiv R(0)=0$, so $L\equiv0$, so
$S'=AS$ with $S(0)=I$, and $S=P$.
If $E=h^{p}K+O(h^{p+1})$ with $K\ne0$, then \eqref{eq:LE} gives
$L=p\,h^{p-1}K+O(h^{p})$, and $R$ has the same leading term. Since
$R(0)=0$, this forces $p\ge2$, and then
$R'=p(p-1)h^{p-2}K+O(h^{p-1})$ vanishes to order exactly $p-2$.

\emph{Step 3: rewriting in terms of $D$.}
Since $C=(D-B)/2$, we have $\ad_C=\tfrac12(Z-T)$ and
$Q=\tfrac12[B,D]=\tfrac12TD$. Hence
$$
R'(h)=\sum_{j\ge0}\frac{h^{j}}{j!\,2^{\,j+2}}\,(b_j-a_j),
\qquad
a_j=(Z-T)^{j}TD,\quad b_j=(-T)^{j}TD .
$$
So the splitting is exact iff $a_j=b_j$ for all $j$, and otherwise
$p-2=\min\{j:a_j\ne b_j\}$.

\emph{Step 4: induction.}
Suppose $a_{j-1}=b_{j-1}$ for some $j\ge1$. Then
$$
a_j=(Z-T)\,b_{j-1}=(-1)^{j-1}Z\,T^{j}D+b_j .
$$
As $a_0=b_0$, it follows that $a_j=b_j$ for all $j<s^*$, where
$s^*=\min\{s\ge1:ZT^{s}D\ne0\}$, and $a_{s^*}\ne b_{s^*}$ if $s^*$ is
finite. Thus $p-2=s^*$, and the splitting is exact iff $s^*=\infty$.
The last claim is the case $s=1$: $ZTD=[D,[B,D]]$.
\end{proof}

\begin{remark}
Since $[D,e^{hB}De^{-hB}]=\sum_{s\ge1}\frac{h^{s}}{s!}ZT^{s}D$, the
criterion reads
$$
p(B,C)=2+\operatorname{ord}_{h=0}\,\bigl[D,\;e^{hB}De^{-hB}\bigr],
$$
and the splitting is exact iff $D$ commutes with all its conjugates
$e^{hB}De^{-hB}$. The leading error coefficient for $p=3$ is
$K=-\tfrac1{48}[D,[B,D]]=-\tfrac1{24}[B+2C,[B,C]]$.
\end{remark}

\section{Upper bound}

\begin{lemma}\label{lem:arb}
For $m\times m$ matrices with nonexact splitting, $p(B,C)\le m(m-1)+1$.
\end{lemma}

\begin{proof}
Let $W_j=\spn\{T^{s}D:1\le s\le j\}$ and $W=\bigcup_jW_j$. If
$T^{j+1}D\in W_j$ then $TW_j\subseteq W_j$ and $W=W_j$. Hence
$\dim W_j=j$ until the chain stabilises, and $W=W_\ell$ with
$\ell=\dim W$.

Every $T^{s}D$ with $s\ge1$ lies in $\im T$, so
$\ell\le\operatorname{rank}\ad_B=m^{2}-\dim\ker\ad_B\le m^{2}-m$.
The last inequality holds because $\dim\ker\ad_B=m$ when $B$ has
distinct eigenvalues (the centraliser is then the algebra of polynomials
in $B$), such $B$ are dense, and $\operatorname{rank}\ad_B$ is lower
semicontinuous in $B$.

If $ZT^{s}D=0$ for $1\le s\le\ell$, then $Z$ vanishes on $W_\ell=W$ and
the splitting is exact by Lemma~\ref{lem:crit}. So a nonexact pair has
$s^*\le\ell\le m(m-1)$, i.e.\ $p\le m(m-1)+2$. This number is even and
$p$ is odd (Lemma~\ref{lem:parity}).
\end{proof}

\begin{proof}[Proof of the upper bound in Theorem~\ref{thm:main}]
Let $B,C$ be $n\times n$ Markov generators with nonexact splitting and
$p>3$. By Lemma~\ref{lem:crit}, $Z^{2}B=-[D,[B,D]]=0$.

\emph{$B$ and $C$ preserve the generalised eigenspaces of $D$.}
Write $\mathbb C^{n}=\bigoplus_\lambda E_\lambda$ with $E_\lambda$ the
generalised eigenspaces of $D$, and let
$B_{\lambda\mu}:E_\mu\to E_\lambda$ be the blocks of $B$. With
$D|_{E_\lambda}=\lambda+N_\lambda$, $N_\lambda$ nilpotent,
$$
(ZB)_{\lambda\mu}=\Phi(B_{\lambda\mu}),\qquad
\Phi(X)=(\lambda-\mu)X+(N_\lambda X-XN_\mu).
$$
Left multiplication by $N_\lambda$ and right multiplication by $N_\mu$
are commuting nilpotent operators, so $\Phi$ is invertible when
$\lambda\ne\mu$. Then $\Phi^{2}(B_{\lambda\mu})=0$ forces
$B_{\lambda\mu}=0$. So $B$, and hence $C=(D-B)/2$, is block diagonal.

\emph{The block at $0$ is exact.}
$D=B+2C$ is a Markov generator, so $e^{tD}$ is a stochastic matrix for
$t\ge0$ and $\|e^{tD}\|_\infty=1$. A Jordan block of size $\ge2$ at the
eigenvalue $0$ would make $e^{tD}$ grow linearly in $t$. Hence
$E_0=\ker D$, and $E_0\ni\one$ is nonzero. On $E_0$ we have
$C=-B/2$, so $B$ and $C$ commute there and $S=P$ on $E_0$.

\emph{Conclusion.}
Let $W'=\bigoplus_{\lambda\ne0}E_\lambda$, of dimension
$m\le n-1$. Since $S-P=0\oplus(S-P)|_{W'}$, the restricted pair is
nonexact with the same order, and Lemma~\ref{lem:arb} gives
$p\le m(m-1)+1\le(n-1)(n-2)+1$.
\end{proof}

\section{Lower bound}

\begin{lemma}[Shift--corner pair]\label{lem:shift}
Let $m\ge2$, let $E_{ij}$ be the matrix units, and set
$$
B=F=\sum_{i=1}^{m-1}E_{i+1,i},\qquad D=E_{1m},\qquad C=\tfrac12(D-F).
$$
Then $B+2C=D$, the splitting is nonexact, and
$$
p(B,C)=\begin{cases}m+1,&m\text{ even},\\ m+2,&m\text{ odd}.\end{cases}
$$
\end{lemma}

\begin{proof}
Left and right multiplication by $F$ commute, and
$F^{a}E_{1m}F^{b}=E_{1+a,\,m-b}$ for $0\le a,b\le m-1$ (zero otherwise).
Hence
$$
T^{s}D=\sum_{\substack{a+b=s\\0\le a,b\le m-1}}(-1)^{b}\binom sb\,
E_{1+a,\,m-b}.
$$
Also $ZE_{ij}=\delta_{mi}E_{1j}-\delta_{j1}E_{im}$, so $Z$ annihilates
$E_{1+a,m-b}$ unless $a=m-1$ or $b=m-1$. Therefore:
\begin{itemize}
\item $ZT^{s}D=0$ for $1\le s\le m-2$;
\item for $s=m-1$ the terms $(a,b)=(m-1,0)$ and $(0,m-1)$ give
$ZT^{m-1}D=\bigl(1+(-1)^{m}\bigr)D$;
\item for $s=m$ and $m\ge 3$ the terms $(m-1,1)$ and $(1,m-1)$ give
$ZT^{m}D=-m\,E_{1,m-1}+(-1)^{m}m\,E_{2,m}$.
\end{itemize}
If $m$ is even, the first nonzero term is at $s=m-1$. If $m$ is odd, it
is at $s=m$. Apply Lemma~\ref{lem:crit}.
\end{proof}

\begin{lemma}[Markov embedding]\label{lem:embed}
Let $B,C$ be real $m\times m$ matrices and $n=m+1$. There are
$n\times n$ Markov generators $B',C'$ with strictly positive
off-diagonal entries whose splitting is exact iff that of $B,C$ is, and
has the same order otherwise.
\end{lemma}

\begin{proof}
Let $U$ be an $n\times m$ matrix whose columns are an orthonormal basis
of $\one^\perp$, and $J=\frac1n\one\one^{\top}$. Then
$U^{\top}U=I_m$, $U^{\top}\one=0$, $UU^{\top}=I_n-J$. For $\lambda>0$ put
$$
B'=U(B-\lambda I_m)U^{\top}=UBU^{\top}+\lambda(J-I_n),\qquad
C'=U(C-\lambda I_m)U^{\top}.
$$
Row sums vanish because $U^{\top}\one=0$. Each off-diagonal entry is an
entry of $UBU^{\top}$ or $UCU^{\top}$ plus $\lambda/n$, hence positive
for $\lambda$ large.

Since $(UXU^{\top})(UYU^{\top})=UXYU^{\top}$, we get
$e^{UXU^{\top}}=J+Ue^{X}U^{\top}$, and $JU=0$. Scalars commute with
everything, so
$$
e^{hB'/2}e^{hC'}e^{hB'/2}-e^{h(B'+C')}
=e^{-2\lambda h}\,U\bigl(S(h)-P(h)\bigr)U^{\top}.
$$
The factor $e^{-2\lambda h}$ never vanishes and $X\mapsto UXU^{\top}$ is
injective, so exactness and the order of vanishing are preserved.
\end{proof}

\begin{proof}[Proof of the lower bound and of $M_2=M_3=3$]
For $n\ge3$, apply Lemma~\ref{lem:embed} to the pair of
Lemma~\ref{lem:shift} with $m=n-1$. The order is $n$ if $n$ is odd and
$n+1$ if $n$ is even, i.e.\ $2\lfloor n/2\rfloor+1$.

For $n=2$ take
$$
B=\begin{pmatrix}-1&1\\1&-1\end{pmatrix},\qquad
C=\begin{pmatrix}-1&1\\2&-2\end{pmatrix},\qquad
[D,[B,D]]=16\begin{pmatrix}1&-1\\1&-1\end{pmatrix}\ne0,
$$
so $p=3$ by Lemma~\ref{lem:crit}.

For $n=2,3$ both bounds equal $3$.
\end{proof}

\end{document}
